\section{Introduction}
\label{sec:intro}

In electric power systems and wholesale electricity markets, the operational value of wind power forecasting is governed by asymmetric decision costs arising from reserve allocation, economic dispatch commitments, and real-time balancing settlement \cite{holttinen2021,weron2014}. Under standard market clearing and system balancing rules, reserve shortfall penalties for under-delivery are typically substantially more severe than downward redispatch or spill penalties for over-delivery \cite{zhang2014wind}. In these asymmetric operational regimes, conventional symmetric regression metrics, such as Root Mean Squared Error (RMSE) and Mean Absolute Error (MAE), fail to reflect tail operational risk and downside financial exposure \cite{lipton2019,pinson2007trading}. To improve generalization and enforce physically realistic outputs, an expanding body of literature embeds aerodynamic limits, empirical power curves, kinematic wake equations, and joint probability distribution constraints directly into deep neural network architectures and loss functions \cite{bastankhah2014new,karniadakis2021physics,raissi2019physics,s34_prx2023,pct2026}. A foundational premise underpins this research paradigm: that reported reductions in forecasting error can be directly attributed to the valid physical domain knowledge encoded within the constraint.

However, observing that a physics-constrained model outperforms an unconstrained baseline does not establish that the gain arises from genuine domain knowledge. In empirical machine learning, performance differences between candidate models can be driven by multiple competing mechanisms:
\begin{enumerate}[label=(\alph*),leftmargin=*]
\item \textit{Parametric capacity shifts}: modifying neural submodules or layer architectures when adding constraints confounds physical bias with representational capacity;
\item \textit{Information-budget leakage}: fitting empirical reference priors on historical datasets provides extra observation volume unavailable to the base forecaster;
\item \textit{Non-specific numerical regularization}: additional penalty objectives can act as generic smoothers or L2 regularizers rather than imparting domain aerodynamic understanding;
\item \textit{Structural approximation error}: fixed algebraic equations or static reference curves inevitably introduce approximation bias as data availability grows;
\item \textit{Temporal proxy misalignment}: regularizing offline training targets against lagged weather telemetry introduces state distortion when forward meteorological forecasts are absent;
\item \textit{Decision-metric misalignment}: improvements under symmetric point metrics frequently fail to translate into downside risk reductions under asymmetric market dispatch.
\end{enumerate}
Because conventional benchmarking protocols rarely isolate these confounding vectors, existing literature leaves the foundational attribution question unresolved: When does a physical prior provide genuine decision value, when is its apparent benefit non-specific, and when does it become harmful?

In this investigation, we address this question by evaluating physical priors as \textbf{conditional inductive guardrails} rather than universally beneficial performance enhancers. The central scientific problem is to identify when a physical prior provides genuine decision value, when its apparent benefit is non-specific, and when it becomes harmful. To delineate this operational envelope, we introduce a \textbf{Controlled Attribution Protocol} across three commercial UK wind farms (Kelmarsh, Hill of Towie, and Penmanshiel) coupled with 2017 UK Elexon Balancing and Settlement Code (BSC) market clearing prices. Holding forecasting backbones, neural capacity, data splits, and post-hoc quantile calibration strictly constant, we contrast authentic empirical power curve priors against form- and weight-matched pseudo-scrambled negative controls across varied sample regimes, constraint weights, and proxy delays.

Our controlled experiments establish that physical priors exhibit conditional, context-dependent operational utility:
\begin{enumerate}[leftmargin=*]
\item \textbf{Prior Specificity under Sample Restriction}: In the tested sample-restricted regime on Penmanshiel ($N_{\text{train}}=1{,}583$), strong empirical curve constraints ($\lambda=1.0$) reduce normalized asymmetric dispatch loss by 1.28\%, whereas permuting the aerodynamic mapping via a pseudo-scrambled control inflates dispatch loss by 61.2\%, providing empirical evidence of prior-specific mapping effects in this sample-restricted setting.
\item \textbf{Diminishing Advantages across Sample Budgets}: Across within-site learning curves ($N_{\text{train}} \in [500, 5000]$), fixed empirical priors provide cost reductions in the tested restricted windows ($N_{\text{train}} \in \{500, 1{,}000\}$). A strictly budget-matched information control ($N_{\text{prior}}=N_{\text{train}}=500$) confirms that domain-structured prior effects retain 77\% to 92\% of the full-history benefit, demonstrating that the majority of the observed benefit cannot be explained by additional historical data volume alone. However, at the largest evaluated training budget ($N_{\text{train}}=5{,}000$), unconstrained data-driven models match or slightly outperform fixed priors (+0.33\% on Kelmarsh, +1.73\% on Hill of Towie), consistent with a transition where the regularizing value of a fixed prior diminishes while structural approximation error becomes increasingly consequential.
\item \textbf{Non-Specific Numerical Regularization Mimicry}: Under weak penalty scaling ($\lambda=0.01$), authentic and scrambled priors yield closely matched outcomes ($0.08803$ vs.\ $0.08792$ on Kelmarsh; $0.08913$ vs.\ $0.08910$ on Hill of Towie), indicating that the observed effect cannot be specifically attributed to physical mapping and remains consistent with non-specific numerical regularization.
\item \textbf{Empirical Boundary Crossovers}: Continuous stress testing on Penmanshiel reveals that the evaluated corruption trajectory crosses the unconstrained baseline around an empirical crossover region of $\alpha \approx 0.07$, while turbine capacity derating exceeding 10\% flips physical priors into net economic penalties under the tested configuration. In the Hill of Towie lag scan, decision risk exceeds the unconstrained reference by the first evaluated lag ($k=6$ / 1\,h delay) and deteriorates monotonically through $k=36$.
\item \textbf{Prediction--Decision Decoupling}: In an info-symmetric evaluation on Kelmarsh, adding wind information improves RMSE by 1.22\% yet worsens asymmetric decision cost by 1.35\%. Downstream, a degenerate zero-bid baseline ($q \equiv 0$) mechanically generates six-figure apparent balancing cashflows without forecasting skill, exposing standalone unhedged balancing cashflows as an unreliable proxy for operational performance.
\end{enumerate}

To contextualize these empirical failure regimes within current reporting practices, we conduct an evidence-gap literature audit of 46 recent studies (2023--2026) and a supporting static computational graph audit of the sole open-source end-to-end training codebase in this domain. Studies in this scoping cohort reported exclusively within nominal operating envelopes (0/46 documenting operational failure regimes), while 44/46 lacked public code and 50.0\% relied exclusively on confounded module ablations. Concurrently, static AST inspection of the Physics-Constrained Transformer (PCT) uncovers an inverted argument order in its divergence loss, producing an analytical $\mathcal{O}(1/\epsilon)$ gradient response under small predicted probabilities while its bivariate kernel coupling remains mathematically sound.

\subsection{Summary of Contributions}
This investigation establishes an evidence-calibrated foundation for physics-constrained wind power forecasting through three interconnected contributions:
\begin{itemize}[leftmargin=*]
\item \textbf{Controlled Attribution Protocol for Physical Priors}: We formulate and execute a controlled benchmarking framework that pairs authentic physical priors with form- and weight-matched pseudo-scrambled negative controls under identical network architectures, optimization schedules, and post-hoc quantile calibrations, providing an empirical standard to distinguish domain aerodynamic guidance from non-specific numerical smoothing.
\item \textbf{Empirical Operating Envelope Mapping}: Across three commercial wind farms, we map the operational boundaries of empirical power-curve priors: documenting the attenuation of prior advantages as sample size increases, confirming via 500-sample budget matching that benefits cannot be explained by historical volume alone, and characterizing continuous prior corruption ($\alpha \approx 0.07$ crossover on Penmanshiel) and training proxy staleness (where decision risk exceeds the unconstrained reference by the first evaluated lag $k=6$ / 1\,h delay on Hill of Towie and deteriorates monotonically through $k=36$).
\item \textbf{Decision-Aligned Evaluation and Governance Framework}: Integrating a decision-theoretic risk decomposition, formal counterexamples, and empirical evaluations (info-symmetric decoupling and zero-bid reductio), we demonstrate why symmetric point metrics and standalone balancing cashflows fail to reflect asymmetric dispatch costs. We translate these findings into a five-step practical decision rule and a five-dimension methodological reporting checklist (C1--C5) for physics-constrained forecasting.
\end{itemize}

\paragraph{Empirical Scope and Methodological Boundaries}
To prevent over-generalization, we explicitly delineate four empirical boundaries governing this study: first, our literature audit characterizes reporting patterns within the 46 studies identified under our prespecified PRISMA-ScR protocol (2023--2026), rather than asserting universal claims across all scientific machine learning; second, software auditing is focused on the computational graph of the open-source PCT benchmark; third, prospective empirical ablations evaluate boundary behaviors across three commercial UK wind farms under matched DLinear backbones, with PatchTST serving as an architectural sensitivity probe; and fourth, downstream evaluations analyze operational decision risk and stylized imbalance cashflow component diagnostics under UK BSC prices, characterizing metric decoupling rather than full commercial portfolio trading.
